\documentclass[../../main.tex]{subfiles}

\begin{document}

\chapter[Context-Free Grammar]{Context-Free Grammar: BNF, EBNF,\\ Derivasi, dan Parse Tree}
\label{chap:cfg}

\begin{outcomebox}
\textbf{Sub-CPMK:} \subcpmk{3.1}, \subcpmk{3.4}\\[3pt]
Setelah menyelesaikan bab ini, mahasiswa mampu:
\begin{enumerate}[label=\arabic*.,leftmargin=*]
  \item mendefinisikan CFG dan menulis tata bahasa dalam BNF;
  \item mengonversi EBNF menjadi BNF dan menyimpan berkas grammar;
  \item menghasilkan derivasi leftmost dan rightmost;
  \item membangun dan menguji pohon sintaks yang dihasilkan;
  \item membuat berkas grammar bebas konteks untuk deklarasi,
        ekspresi, dan pernyataan Subset-C.
\end{enumerate}
\end{outcomebox}

\section*{Materi Pokok}
\begin{itemize}[leftmargin=*]
  \item Definisi CFG, nonterminal, terminal, dan aturan produksi.
  \item Notasi BNF dan EBNF beserta konversinya.
  \item Derivasi leftmost, rightmost, dan pohon urai.
  \item Pengujian derivasi dan pohon sintaks.
  \item Berkas grammar Subset-C untuk deklarasi, ekspresi, dan
        pernyataan.
\end{itemize}

\subfile{section-01}
\subfile{section-02}
\subfile{section-03}
\subfile{section-04}
\subfile{section-05}

\section*{Contoh}
\addcontentsline{toc}{section}{Contoh}

\begin{contoh}[Derivasi dengan tanda kurung]
Untuk string \((\code{id} + \code{id}) * \code{id}\), derivasi
leftmost pada tata bahasa aritmetika adalah:
\[
\begin{aligned}
E &\Rightarrow T\\
  &\Rightarrow T * F\\
  &\Rightarrow F * F\\
  &\Rightarrow (E) * F\\
  &\Rightarrow (E + T) * F\\
  &\Rightarrow (T + T) * F\\
  &\Rightarrow (F + T) * F\\
  &\Rightarrow (\code{id} + T) * F\\
  &\Rightarrow (\code{id} + F) * F\\
  &\Rightarrow (\code{id} + \code{id}) * F\\
  &\Rightarrow (\code{id} + \code{id}) * \code{id}.
\end{aligned}
\]
\end{contoh}

\begin{contoh}[Pohon urai if-else]
Diberikan pernyataan \code{if (x) y; else z;}. Pohon urai
tingkat atas memiliki simpul \code{ifstmt} dengan anak
\code{"if"}, \code{"("}, ekspresi \code{x}, \code{")"},
pernyataan \code{y;}, dan bagian \code{else z;}. Jika tata bahasa
tidak membedakan pasangan \code{else} terdekat, pohon alternatif
dapat terbentuk. Ambiguitas ini dibahas pada Bab IX.
\end{contoh}

\begin{contoh}[Konversi EBNF]
Aturan EBNF
\[
\text{args} = \text{expr} \{\text{"," expr}\}
\]
dikonversi menjadi:
\[
\begin{aligned}
\text{args} &\to \text{expr}\ \text{args\_tail},\\
\text{args\_tail} &\to \text{"," expr}\ \text{args\_tail}
\mid \varepsilon.
\end{aligned}
\]
Aturan ini right-recursive, sehingga cocok untuk parser top-down.
\end{contoh}

\begin{contoh}[AST ekspresi]
Untuk \(a + b * c\), AST memiliki akar \code{+}, anak kiri
\code{a}, dan anak kanan \code{*} dengan anak \code{b} dan
\code{c}. Jika ditulis dalam notasi prefix:
\[
\code{(+ a (* b c))}.
\]
Notasi ini menghilangkan tanda kurung sintaksis dan langsung
menunjukkan struktur operator.
\end{contoh}

\begin{contoh}[Merancang aturan fungsi]
Deklarasi
\begin{lstlisting}[language=C]
int tambah(int a, int b) {
    return a + b;
}
\end{lstlisting}
diurai oleh aturan \code{topdecl}: tipe \code{int}, identifier
\code{tambah}, daftar parameter, lalu blok. Daftar parameter
menggunakan \code{params} dan \code{params\_tail}. Blok berisi satu
pernyataan \code{return}. Setiap bagian memiliki nonterminal sendiri,
sehingga parser dapat memeriksa struktur secara bertahap.
\end{contoh}

\section*{Aktivitas Pembelajaran dan Praktikum}
\addcontentsline{toc}{section}{Aktivitas Pembelajaran dan Praktikum}

\begin{table}[H]
\centering
\caption{Rencana aktivitas Bab VIII}
\label{tab:ch08-aktivitas}
\begin{tabularx}{\textwidth}{@{}llX@{}}
\toprule
\textbf{Sesi} & \textbf{Bentuk} & \textbf{Kegiatan} \\
\midrule
1 & Ceramah & CFG, BNF, dan derivasi \\
2 & Latihan & Menurunkan string dan menggambar pohon urai \\
3 & Lokakarya & Konversi EBNF ke BNF \\
4 & Praktikum & P2-A: tata bahasa Subset-C dan peer review \\
\bottomrule
\end{tabularx}
\end{table}

\subsection*{P2-A: Tata Bahasa Subset-C}
\textbf{Tujuan.} Menulis tata bahasa EBNF untuk Subset-C, lalu
mengonversinya ke BNF.

\textbf{Langkah kerja.}
\begin{enumerate}[label=\arabic*.,leftmargin=*]
  \item Lengkapi \file{spec/grammar.ebnf} dengan aturan untuk seluruh
        konstruksi yang didukung.
  \item Turunkan lima string contoh: ekspresi, deklarasi, if, while,
        dan fungsi.
  \item Konversi EBNF ke BNF pada \file{spec/grammar.bnf}.
  \item Gambar pohon urai untuk dua string.
  \item Tukar berkas dengan teman untuk peer review.
\end{enumerate}

\textbf{Artefak.} Berkas \file{spec/grammar.ebnf},
\file{spec/grammar.bnf}, lima derivasi, dua pohon urai, dan catatan
hasil uji.

\textbf{Definisi selesai.}
\begin{checklist}
  \item EBNF mencakup semua konstruksi: deklarasi, ekspresi, dan
        pernyataan.
  \item Lima derivasi benar dan terdokumentasi.
  \item BNF konsisten dengan EBNF.
  \item Dua pohon urai benar dan sesuai derivasi.
  \item Berkas grammar dan bukti pengujian dikumpulkan.
  \item Catatan peer review lengkap.
\end{checklist}

P2-A menjadi bagian Praktikum 2 (12\%) yang dilaporkan pada Bab X.

\section*{Latihan dan Refleksi}
\addcontentsline{toc}{section}{Latihan dan Refleksi}

\subsection*{Latihan}
\begin{enumerate}[label=\arabic*.,leftmargin=*]
  \item Definisikan CFG dan jelaskan perbedaan nonterminal dan
        terminal.
  \item Tuliskan BNF untuk bahasa \(\{a^n b^n \mid n \ge 0\}\).
  \item Konversi aturan EBNF \([X]\ Y\) ke BNF.
  \item Konversi aturan EBNF \(\{X\}\ Y\) ke BNF.
  \item Buat derivasi leftmost dan rightmost untuk
        \((\code{id} + \code{id}) * \code{id}\).
  \item Gambarkan pohon urai untuk string pada Latihan 5.
  \item Jelaskan perbedaan pohon urai dan AST.
  \item Rancang aturan untuk pernyataan \code{for}.
\end{enumerate}

\subsection*{Refleksi}
\begin{reflectionbox}
\begin{enumerate}[label=\arabic*.,leftmargin=*]
  \item Mengapa tata bahasa dengan tingkat prioritas terpisah lebih
        mudah diurai?
  \item Apakah Anda dapat menjelaskan mengapa AST lebih kecil daripada
        pohon urai?
  \item Bagian mana dari perancangan tata bahasa yang paling menantang?
\end{enumerate}
\end{reflectionbox}

\section*{Asesmen}
\addcontentsline{toc}{section}{Asesmen}

Asesmen Bab VIII mencakup Praktikum 2 (12\%). Fokusnya adalah
ketepatan berkas grammar, derivasi, dan pohon sintaks.

\begin{table}[H]
\centering
\caption{Instrumen asesmen Bab VIII}
\label{tab:ch08-asesmen}
\begin{tabularx}{\textwidth}{@{}lXl@{}}
\toprule
\textbf{Instrumen} & \textbf{Indikator} & \textbf{Bobot} \\
\midrule
Praktikum 2 & Berkas EBNF/BNF, lima derivasi, dua pohon urai, dan
catatan uji & 12\% (bagian) \\
Kuis Konsep & Definisi CFG dan pohon urai & formatif \\
\bottomrule
\end{tabularx}
\end{table}

Kriteria kelulusan: tata bahasa mencakup seluruh konstruksi,
derivasi benar, pohon sintaks konsisten dengan tata bahasa, dan bukti
pengujian dikumpulkan.

\section*{Checklist Pencapaian Kompetensi}
\addcontentsline{toc}{section}{Checklist Pencapaian Kompetensi}

\begin{checklist}
  \item Saya dapat mendefinisikan CFG.
  \item Saya dapat menulis aturan dalam BNF.
  \item Saya dapat membaca dan menulis EBNF.
  \item Saya dapat mengonversi EBNF ke BNF.
  \item Saya dapat menghasilkan derivasi leftmost dan rightmost.
  \item Saya dapat menggambar pohon urai.
  \item Saya dapat menjelaskan perbedaan pohon urai dan AST.
  \item Saya dapat merancang tata bahasa untuk Subset-C.
  \item Saya sudah menyelesaikan P2-A.
\end{checklist}

\section*{Rangkuman}
\addcontentsline{toc}{section}{Rangkuman}

Pada bab ini telah dibahas:
\begin{itemize}[leftmargin=*]
  \item definisi CFG dan notasi BNF;
  \item EBNF dan aturan konversinya ke BNF;
  \item derivasi leftmost, rightmost, dan bentuk kalimat;
  \item pohon urai dan hubungannya dengan AST;
  \item perancangan tata bahasa Subset-C.
\end{itemize}

Bab IX akan membersihkan tata bahasa dari ambiguitas, rekursi kiri,
dan faktor umum. Transformasi ini diperlukan agar tata bahasa dapat
diurai secara top-down pada Bab X.

\section*{Evaluasi}
\addcontentsline{toc}{section}{Evaluasi}

\begin{enumerate}[label=\arabic*.,leftmargin=*]
  \item Definisikan CFG dan berikan satu contoh.
  \item Tuliskan BNF untuk ekspresi aritmetika dengan prioritas.
  \item Konversi aturan EBNF \code{params = param \{"," param\}} ke
        BNF.
  \item Buat derivasi leftmost untuk \code{id * (id + id)}.
  \item Gambarkan pohon urai untuk string pada Soal 4.
  \item Jelaskan perbedaan pohon urai dan AST dengan contoh.
  \item Rancang aturan untuk deklarasi variabel dengan inisialisasi
        opsional.
  \item Mengapa tata bahasa Subset-C perlu diubah sebelum parsing
        top-down?
  \item Buat berkas grammar EBNF/BNF untuk deklarasi, ekspresi, dan
        pernyataan; tunjukkan bukti derivasi dan pohon untuk dua
        string.
\end{enumerate}

\subsection*{Kunci dan Rubrik Singkat}
\begin{itemize}[leftmargin=*]
  \item Soal 1--3: ketepatan definisi dan konversi (70\%),
        kelengkapan (30\%).
  \item Soal 4--5: kebenaran derivasi dan pohon (70\%),
        konsistensi (30\%).
  \item Soal 6--9: ketepatan konsep dan artefak (60\%), contoh
        (40\%).
\end{itemize}


\ifSubfilesClassLoaded{%
  \bibliographystyle{apalike}%
  \bibliography{../../references}%
}{}

\end{document}
